\chapter{Conclusiones}

\begin{enumerate}
    \item El sistema base de interacción en realidad virtual 
    (VR, por sus siglas en inglés) permite reconocer y registrar acciones
    diferenciandas mediante controladores, colisiones y reglas específicas
    para las mecánicas implementadas.

    \item El mecanismo de telemetría registra 5048 interacciones de 34 
    participantes, incluyendo aciertos, fallos, omisiones y métricas 
    específicas de tiempo, reacción, dirección y uso de mano correcta
    para minijuegos correspondientes.

    \item El modelado del jugador representa variaciones del comportamiento
    mediante los estados \textit{idle}, \textit{engaged} y \textit{overloaded},
    generando información operativa para las decisiones del sistema 
    adaptativo de dificultad dinámica (DDA, por sus siglas en inglés).

    \item El DDA realiza 216 modificaciones de dificultad y produce cambios
    observables en el desempeño, con una mejora promedio de 27.35 puntos
    porcentuales después de las reducciones analizables y una disminución 
    de 31.30 puntos después de los aumentos. 

    \item El sistema implementado completa el ciclo de medición, 
    interpretación y adaptación, aunque los resultados subjetivos de
    atención, diversión y agobio no pueden atribuirse causalmente al DDA. 
\end{enumerate}